\documentclass[11pt,letterpaper]{article}
% preamble.tex — Shared preamble for all documents
\usepackage{graphicx}
\usepackage{hyperref}
\usepackage{enumitem}
\usepackage{titlesec}
\usepackage{fancyhdr}
\usepackage{setspace}
\usepackage{parskip}
\usepackage{etoolbox}
\usepackage{xspace}
\usepackage{needspace}
\usepackage{changepage}
\usepackage{geometry}
\usepackage{ragged2e}
\usepackage{environ}

% Hyperref setup
\hypersetup{
  colorlinks=false,
  pdfborder={0 0 0},
  pdftitle={The Letter to the Father},
  pdfauthor={},
}

% Load project files
% colors.tex — Color palette for The Letter to the Father
\usepackage{xcolor}

% Primary palette — muted, period-appropriate, prestige
\definecolor{kafkablack}{HTML}{1A1A1A}
\definecolor{kafkagray}{HTML}{4A4A4A}
\definecolor{kafkalight}{HTML}{8A8A8A}
\definecolor{kafkacream}{HTML}{F5F0E8}
\definecolor{kafkasepia}{HTML}{704214}
\definecolor{kafkablood}{HTML}{6B1C23}
\definecolor{kafkaink}{HTML}{2C2C3A}

% Accent colors — for package materials
\definecolor{prestigegold}{HTML}{B8860B}
\definecolor{sectionrule}{HTML}{3A3A3A}
\definecolor{pullquote}{HTML}{5A4A3A}

% metadata.tex — Project metadata for The Letter to the Father
\newcommand{\filmtitle}{The Letter to the Father}
\newcommand{\filmtitleupper}{THE LETTER TO THE FATHER}
\newcommand{\filmsubtitle}{An Original Screenplay}
\newcommand{\filmauthor}{Based on the life and writings of Franz Kafka}
\newcommand{\filmyear}{2026}
\newcommand{\filmdraft}{Revised Draft -- Public Development Example}
\newcommand{\filmdate}{October 2026 -- Example Edition}
\newcommand{\filmlogline}{In the final years of his life, a frail and tormented Franz Kafka struggles to write, love, and exist within a world of family pressure, bureaucratic humiliation, illness, and spiritual dread --- while the stories he creates begin to mirror and consume the life he can no longer control.}
\newcommand{\filmgenre}{Prestige Drama / Literary Biopic}
\newcommand{\filmrating}{R}
\newcommand{\filmlead}{Adrien Brody}
\newcommand{\filmbudgetrange}{\$12--18M}
\newcommand{\filmruntime}{Approximately 110 minutes}


% Public example status travels with every downloadable document.
\AddToHook{shipout/foreground}{\put(40,-776){\makebox(0,0)[l]{\fontsize{8}{9}\selectfont\textit{Development example. All casting, financing and awards language is illustrative; no attachments confirmed.}}}}

% typography.tex — Typography settings
\usepackage{fontspec}
\usepackage{microtype}

% Prefer repo-local metric-compatible fonts so Tectonic does not depend on
% absolute host-system font paths. Fall back to system fonts only if the vendored
% font payload is unavailable.
\IfFileExists{../../../fonts/Tinos-Regular.ttf}{
  \setmainfont{Tinos-Regular.ttf}[
    Path=../../../fonts/,
    BoldFont=Tinos-Bold.ttf,
    ItalicFont=Tinos-Italic.ttf,
    BoldItalicFont=Tinos-BoldItalic.ttf,
  ]
}{
  \IfFontExistsTF{TeX Gyre Termes}{
    \def\projectmainfont{TeX Gyre Termes}
  }{
    \IfFontExistsTF{Times New Roman}{
      \def\projectmainfont{Times New Roman}
    }{
      \def\projectmainfont{Latin Modern Roman}
    }
  }
  \setmainfont{\projectmainfont}
}

\IfFileExists{../../../fonts/Arimo-Regular.ttf}{
  \setsansfont{Arimo-Regular.ttf}[
    Path=../../../fonts/,
    BoldFont=Arimo-Bold.ttf,
    ItalicFont=Arimo-Italic.ttf,
    BoldItalicFont=Arimo-BoldItalic.ttf,
  ]
}{
  \IfFontExistsTF{TeX Gyre Heros}{
    \def\projectsansfont{TeX Gyre Heros}
  }{
    \IfFontExistsTF{Helvetica}{
      \def\projectsansfont{Helvetica}
    }{
      \def\projectsansfont{Latin Modern Sans}
    }
  }
  \setsansfont{\projectsansfont}
}

\IfFileExists{../../../fonts/Cousine-Regular.ttf}{
  \setmonofont{Cousine-Regular.ttf}[
    Path=../../../fonts/,
    BoldFont=Cousine-Bold.ttf,
    ItalicFont=Cousine-Italic.ttf,
    BoldItalicFont=Cousine-BoldItalic.ttf,
    Scale=0.95,
  ]
  \def\screenplayfont{%
    \fontspec{Cousine-Regular.ttf}[%
      Path=../../../fonts/,
      BoldFont=Cousine-Bold.ttf,
      ItalicFont=Cousine-Italic.ttf,
      BoldItalicFont=Cousine-BoldItalic.ttf,
      Scale=1.0,
    ]%
  }
}{
  \IfFontExistsTF{Courier New}{
    \def\projectmonofont{Courier New}
  }{
    \IfFontExistsTF{Courier}{
      \def\projectmonofont{Courier}
    }{
      \def\projectmonofont{Latin Modern Mono}
    }
  }
  \setmonofont{\projectmonofont}[Scale=0.95]
  \def\screenplayfont{\fontspec{\projectmonofont}[Scale=1.0]}
}

% Microtype settings
\microtypesetup{
  protrusion=true,
  expansion=false,
}

% studio_package_macros.tex — Formatting macros for the studio package

% ============================================================
% SECTION STYLING
% ============================================================
\titleformat{\section}
  {\Large\bfseries\color{kafkablack}\sffamily}
  {}
  {0em}
  {}
  [\vspace{2pt}\textcolor{sectionrule}{\rule{\textwidth}{0.5pt}}\vspace{6pt}]

\titleformat{\subsection}
  {\large\bfseries\color{kafkagray}\sffamily}
  {}
  {0em}
  {}

\titleformat{\subsubsection}
  {\normalsize\bfseries\color{kafkagray}\sffamily}
  {}
  {0em}
  {}

% ============================================================
% PULL QUOTE
% ============================================================
% Usage: \pullquote{Text here}
\NewEnviron{pullquote}{%
  \vspace{8pt}%
  \begin{adjustwidth}{0.5in}{0.5in}%
    \color{pullquote}\itshape\large%
    \BODY%
  \end{adjustwidth}%
  \vspace{8pt}%
}

% ============================================================
% HIGHLIGHT BOX
% ============================================================
\newcommand{\highlightbox}[1]{%
  \par\vspace{6pt}%
  \noindent\colorbox{kafkacream}{%
    \parbox{\dimexpr\textwidth-2\fboxsep}{%
      \vspace{4pt}#1\vspace{4pt}%
    }%
  }%
  \par\vspace{6pt}%
}

% ============================================================
% TARGET ENTRY (for director/producer lists)
% ============================================================
% Usage: \target{Name}{Why they fit}{Tier}
\newcommand{\target}[3]{%
  \needspace{4\baselineskip}%
  \subsubsection*{#1 \hfill \textnormal{\small\textit{#3}}}%
  #2\par\vspace{6pt}%
}

% ============================================================
% SCENE CARD (for showcase section)
% ============================================================
% Usage: \scenecard{Title}{Description}
\newcommand{\scenecard}[2]{%
  \needspace{5\baselineskip}%
  \vspace{4pt}%
  \noindent\textbf{\sffamily #1}\\[3pt]%
  #2\par\vspace{8pt}%
  \textcolor{sectionrule}{\rule{0.3\textwidth}{0.25pt}}%
  \vspace{8pt}%
}

% ============================================================
% TAGLINE
% ============================================================
\newcommand{\tagline}[1]{%
  \vspace{4pt}%
  \begin{center}%
    {\large\itshape\color{kafkagray} ``#1''}%
  \end{center}%
  \vspace{4pt}%
}

% ============================================================
% COMP FILM ENTRY
% ============================================================
% Usage: \compfilm{Title (Year)}{Relevance}
\newcommand{\compfilm}[2]{%
  \noindent\textbf{#1}\\%
  {\small #2}\par\vspace{6pt}%
}

% ============================================================
% RISK ENTRY
% ============================================================
% Usage: \risk{Risk}{Mitigation}
\newcommand{\risk}[2]{%
  \noindent\textbf{Risk:} #1\\%
  \textbf{Mitigation:} #2\par\vspace{8pt}%
}

% ============================================================
% PACKAGE GEOMETRY
% ============================================================
\newcommand{\packagegeometry}{%
  \newgeometry{
    left=1in,
    right=1in,
    top=1in,
    bottom=1in,
  }%
}

\geometry{left=0.75in, right=0.75in, top=0.75in, bottom=0.75in}
\pagestyle{empty}
\setlength{\parindent}{0pt}
\setlength{\parskip}{5pt}

\begin{document}

\begin{center}
{\Large\bfseries\sffamily \filmtitleupper}\\[3pt]
{\normalsize\sffamily\color{kafkagray} Sales \& Distribution Brief}
\end{center}
\vspace{2pt}
\textcolor{sectionrule}{\rule{\textwidth}{0.5pt}}
\vspace{4pt}

\textbf{The Buy}\\
A prestige feature built around a clean acquisition proposition: \textbf{Adrien Brody as Franz Kafka} in a role with obvious awards heat, a culturally legible title, and stronger propulsion than the usual literary-biopic lane. The film is structured in six movements but sells simply: a brilliant son can describe every system except the one that formed him. The recurring office case, the undelivered father letter, and the late Berlin movement with Dora give buyers a movie that is marketable as drama, not homework.

\vspace{2pt}

\textbf{Why It Is More Accessible Than the Subject Suggests}\\
This is not a museum-piece writer portrait. The office material gives the first half procedural traction. The Felice/Milena/Dora architecture gives the middle and back half emotional variety. The father-letter engine gives the campaign a memorable trailer spine and title hook. Even for audiences who have never read Kafka, the movie plays as pressure, shame, love, illness, and the cost of remaining undelivered to the people closest to you.

\vspace{2pt}

\textbf{Acquisition Thesis}\\
Brody provides the package's immediate value and campaign logic. Kafka provides worldwide cultural recognition. The script's differentiator is that it turns those assets into a tighter dramatic object than most prestige-biographical comps: a role with six major clip scenes, a rigorous formal concept, and a final movement that lands with warmth before collapse. It fits the exact space where specialty distributors, prestige streamers, and international sales agents still want films that can open on reviews, festivals, performance, and cultural conversation.

\vspace{2pt}

\textbf{Launch Path}\\
Primary launch thesis: Venice Competition, then Telluride/TIFF, then fall platform release. Cannes is viable if a director from the current target tier lands. The movie wants discovery through curation, then a campaign framed around Brody, the role's physical and emotional demands, and the unusual clarity of the father-letter premise.

\vspace{2pt}

\textbf{Audience \& Territory Fit}\\
Core audience: upscale arthouse, prestige-adult drama, awards-season viewers, literary/cultural audience, and Brody-driven specialty crowd. Strongest international territories are German-speaking Europe, France, Italy, the UK, and urban North America. The package also travels well to prestige streaming because the concept can be summarized in one sentence and the campaign story is instantly legible.

\vspace{2pt}

\textbf{Marketing Hooks}\\
\textit{Campaign line}: the definitive Kafka performance. --- \textit{Trailer engine}: the man who could write the Letter to His Father but never deliver it. --- \textit{Visuals}: the office machine, tribunal-like family scenes, the returned Letter in Berlin, Brod's ``burn everything'' confrontation, the final handwritten slips. --- \textit{Positioning}: \textit{The Brutalist} credibility with a cleaner emotional hook and a more classical prestige-release shape.

\vspace{2pt}

\textbf{Comparable Landing Zone}\\
Searchlight, Focus, Sony Pictures Classics, MUBI, or a top-tier streamer awards play. The film can support either a disciplined platform theatrical campaign or a premium acquisition with festival-and-awards value as the primary driver.

\vspace{4pt}
\textcolor{sectionrule}{\rule{\textwidth}{0.5pt}}
\vspace{2pt}
\begin{center}
{\small\color{kafkalight}\textit{A prestige acquisition with a clear title hook, a real performance event, and a cleaner sales story than the genre usually offers.}}
\end{center}

\end{document}
